% (c) Nikita Lisitsa, lisyarus@gmail.com, 2026

\documentclass[10pt]{beamer}

\usepackage[T2A]{fontenc}
\usepackage[russian]{babel}
\usepackage{minted}

\usepackage[most]{tcolorbox}
\tcbuselibrary{minted}

\usepackage{graphicx}
\graphicspath{ {./images/} }

\usepackage{adjustbox}

\usepackage{color}
\usepackage{soul}

\usepackage{hyperref}

\usepackage{amsmath}

\usepackage{tikz}
\usetikzlibrary{decorations}
\usetikzlibrary{decorations.pathreplacing}
\usepackage{xifthen}

\usetheme{metropolis}
\setminted{fontsize=\footnotesize}

\definecolor{red}{rgb}{1,0,0}
\definecolor{green}{rgb}{0,0.5,0}
\definecolor{blue}{rgb}{0,0,1}
\definecolor{magenta}{rgb}{0.75,0,0.75}
\definecolor{codebg}{RGB}{29,35,49}

\makeatletter
\newcommand{\slideimage}[1]{
  \begin{figure}
    \begin{adjustbox}{width=\textwidth, totalheight=\textheight-2\baselineskip-2\baselineskip,keepaspectratio}
      \includegraphics{#1}
    \end{adjustbox}
  \end{figure}
}
\makeatother

\usemintedstyle{tango}

\title{Компьютерная графика}
\subtitle{Лекция 11: HDR, tone mapping, gamma correction, sRGB, color banding, dithering}
\date{2026}

\setbeamertemplate{footline}[frame number]

\usemintedstyle{solarized-light}

\begin{document}

\frame{\titlepage}

\begin{frame}[fragile]
\frametitle{HDR}
\begin{itemize}
\item Проблема: после вычисления освещённости мы можем получить значения компонент цвета большие, чем 1, и при записи в экран/текстуру они будут обрезаны до 1
\pause
\item В реальном мире интенсивность света ничем не ограничена, и между солнечным днём и тёмной ночью может меняться на 15 десятичных порядков
\pause
\item При рендеринге объект может освещаться тусклым ambient освещением, а может сразу десятком ярких источников света
\end{itemize}
\end{frame}

\begin{frame}[fragile]
\frametitle{HDR}
\begin{itemize}
\item High Dynamic Range (HDR) -- термин, означающий большой (не ограниченный \begin{math}[0, 1]\end{math}) диапазон интенсивностей света
\pause
\item Low Dynamic Range (LDR) -- то, чем мы занимались до этого :)
\pause
\item HDR текстура (например, environment map) -- содержит значения, выходящие за диапазон \begin{math}[0, 1]\end{math} (например, 32-bit floating-point, или unsigned normalized + фиксированный множитель)
\pause
\item HDR дисплеи позволяют отобразить более широкий диапазон яркостей (всё равно ограниченный, так что это скорее маркетинговый термин)
\pause
\item HDR рендеринг позволяет отобразить HDR-освещение
\pause
\item Большинство современных графических движков реализуют HDR
\end{itemize}
\end{frame}

\begin{frame}[fragile]
\frametitle{Tone mapping}
\begin{itemize}
\item Чтобы поддержать HDR рендеринг, Нужно превратить диапазон освещённостей \begin{math}[0, \infty)\end{math} в диапазон значений пикселей экрана \begin{math}[0, 1]\end{math}
\pause
\item В принципе, подойдёт любая монотонно возрастающая функция \begin{math}[0, \infty)\rightarrow [0, 1]\end{math} -- \textit{tone-mapping curve}:
\begin{itemize}
\item \begin{math}x \mapsto \frac{x}{1+x}\end{math} -- \textit{Reinhard operator}
\item \begin{math}x \mapsto \arctan(x)\end{math}
\item \begin{math}x \mapsto \frac{2}{1+e^{-x}}-1\end{math}
\end{itemize}
\pause
\item Иногда используют более сложные функции, взятые из кинематографа (filmic tone mapping)
\pause
\item Иногда используют функции, работающие не отдельно на каждом канале, а с цветом целиком
\pause
\item Иногда всё равно ограничивают диапазон входных интенсивностей каким-то максимальным значением \begin{math}W\end{math} (\textit{white}), и делают нелинейное преобразование \begin{math}[0, W] \rightarrow [0, 1]\end{math}
\end{itemize}
\end{frame}

\begin{frame}[fragile]
\frametitle{Reinhard operator}
\slideimage{reinhard.png}
\end{frame}

\begin{frame}[fragile]
\frametitle{Filmic tone mapping curve (ACES)}
\slideimage{aces.png}
\end{frame}

\begin{frame}[fragile]
\frametitle{Uncharted 2 tone mapping}
\usemintedstyle{lightbulb}
\begin{tcblisting}{listing engine=minted, minted language=wgsl, minted options={fontsize=\scriptsize}, listing only, colback=codebg, colframe=codebg, rounded corners}
fn TonemapRaw(x: vec3f) -> vec3f {
    let A = 0.15;
    let B = 0.50;
    let C = 0.10;
    let D = 0.20;
    let E = 0.02;
    let F = 0.30;

    return ((x*(A*x+C*B)+D*E)/(x*(A*x+B)+D*F))-E/F;
}

fn Uncharted2Tonemap(color: vec3f) -> vec3f {
    float W = 11.2;
    return TonemapRaw(color) / TonemapRaw(vec3f(W));
}
\end{tcblisting}
\end{frame}

\begin{frame}[fragile]
\frametitle{Сравнение разных tone-mapping кривых}
\slideimage{tone_mapping.png}
\end{frame}

\begin{frame}[fragile]
\frametitle{Tone mapping}
\begin{itemize}
\item Tone-mapping функции часто позволяют настраивать среднюю/максимальную видимую освещённость сцены либо \textit{выдержку} (\textit{exposure})
\pause
\item Чем выше средняя/максимальная освещённость, тем \textit{темнее} сцена
\pause
\item Чем выше выдержка, тем \textit{светлее} сцена
\pause
\item Выдержку можно подгонять под конкретную сцену, а можно автоматически подстраивать под среднюю яркость видимых камерой объектов (\textit{dynamic exposure, auto-exposure})
\end{itemize}
\end{frame}

\begin{frame}[fragile]
\frametitle{Witcher 3: high exposure}
\slideimage{witcher_hdr_inside.png}
\end{frame}

\begin{frame}[fragile]
\frametitle{Witcher 3: low exposure}
\slideimage{witcher_hdr_outside.png}
\end{frame}

\begin{frame}[fragile]
\frametitle{Как применять tone mapping}
\begin{itemize}
\item При рендеринге без прозрачности tone-mapping можно применять в том же шейдере, который вычисляет освещение, и писать значения цвета сразу на экран
\pause
\item При включенном блендинге (напр. deferred shading с наложением освещения по одному источнику за раз), применять tone-mapping в этом же шейдере нельзя: сумма tone-mapped цветов не равна значению tone map от суммы (т.е. tone-mapping функция нелинейна)
\pause
\item В этом случае обычно результирующий цвет пикселя аккумулируется в HDR-текстуре (\mintinline{rust}|Rgba16Float|, \mintinline{rust}|Rgba32Float|, \mintinline{rust}|Rg11b10Ufloat|, \mintinline{rust}|Rgb9e5Ufloat|), и потом отдельным проходом рисуется на экран с tone-mapping'ом
\pause
\item Некоторые эффекты пост-обработки лучше работают именно в HDR диапазоне (т.е. до применения tone mapping) -- размытие, свечение (bloom), и т.п.
\end{itemize}
\end{frame}

\begin{frame}[fragile]
\frametitle{HDR \& tone mapping: ссылки}
\begin{itemize}
\item \href{https://learnopengl.com/Advanced-Lighting/HDR}{\texttt{learnopengl.com/Advanced-Lighting/HDR}}
\item \href{https://skylum.com/blog/what-is-tone-mapping}{\texttt{What is tone mapping}}
\item \href{https://www.veneratech.com/what-is-hdr-tone-mapping}{\texttt{What is HDR tone mapping}}
\item \href{https://knarkowicz.wordpress.com/2016/01/06/aces-filmic-tone-mapping-curve}{\texttt{ACES filmic tone mapping curve}}
\item \href{http://filmicworlds.com/blog/filmic-tonemapping-operators/}{\texttt{Filmic tone-mapping operators}}
\end{itemize}
\end{frame}

\begin{frame}<1>[fragile,label=gamma]
\frametitle{Гамма}
\begin{itemize}
\item Обычно, интенсивность света \begin{math}I\end{math}, излучаемого монитором, нелинейно зависит от значения \begin{math}V\end{math}, записанного в пикселе
\pause
\item Это лучше соответствует восприятию света человеком и даёт больше точности тёмным цветам, которые человек различает лучше
\pause
\item Почти всегда используется показательная функция:
\begin{equation}I \sim V^\gamma\end{equation}
\pause
\item \begin{math}\gamma\end{math} обычно равна 2.2 (некоторые компьютеры Macintosh использовали 1.8)
\end{itemize}
\end{frame}

\begin{frame}
\frametitle{Линейное значение пикселя vs линейная интенсивность излучения}
\begin{figure}
\slideimage{gamma-scale.png}
\end{figure}
\end{frame}

\againframe<2->{gamma}

\begin{frame}<1>[fragile,label=gamma2]
\frametitle{Проблемы гаммы}
\begin{itemize}
\item Картинка может издалека выглядеть ярче, чем её усреднённый вариант (e.g. mipmap)
\pause
\item Искажается восприятие относительных яркостей, особенно при реалистичном рендеринге (e.g. объект в два раза ярче не будет выглядеть в два раза ярче)
\pause
\item Неправильно выглядит освещение, наложение источников света, и т.д.
\end{itemize}
\end{frame}

\begin{frame}
\frametitle{Серый (цвет=0.5) квадрат и квадрат с мелкой шахматной раскраской}
\begin{figure}
\slideimage{gamma-checkers.png}
\end{figure}
\end{frame}

\againframe<2->{gamma2}

\begin{frame}[fragile]
\frametitle{Коррекция гаммы (gamma-correction)}
\begin{itemize}
\item Коррекция гаммы -- общий термин для применения любых нелинейных преобразований над интенсивностью пикселя
\pause
\item В рендеринге под гамма-коррекцией обычно подразумевают применение обратного к \begin{math}V^{2.2}\end{math} преобразования, чтобы получить линейную зависимость выходящего излучения от значения пикселя
\end{itemize}
\end{frame}

\begin{frame}[fragile]
\frametitle{Коррекция гаммы (gamma-correction)}
\usemintedstyle{lightbulb}
\begin{tcblisting}{listing engine=minted, minted language=wgsl, minted options={fontsize=\scriptsize}, listing only, colback=codebg, colframe=codebg, rounded corners}
// Вычислили цвет пикселя с учётом освещения
var color = ...;

// Гамма-коррекция
color = pow(color, vec3f(1.0 / 2.2));
\end{tcblisting}
\end{frame}

\begin{frame}[fragile]
\frametitle{Эффекты коррекции гаммы}
\slideimage{gamma-ex1.png}
\end{frame}

\begin{frame}[fragile]
\frametitle{Эффекты коррекции гаммы}
\slideimage{gamma-ex2.png}
\end{frame}

\begin{frame}[fragile]
\frametitle{Эффекты коррекции гаммы}
\slideimage{gamma-ex3.png}
\end{frame}

\begin{frame}[fragile]
\frametitle{Эффекты коррекции гаммы}
\begin{itemize}
\item Типичные эффекты гамма-коррекции:
\pause
\begin{itemize}
\item Картинка становится ярче (чем темнее цвет, тем больше прирост яркости)
\pause
\item Картинка становится менее контрастной
\pause
\item Освещение выглядит правильнее
\pause
\item Переход от освещённой к неосвещённой области объекта выглядит более резким (как в реальности)
\end{itemize}
\pause
\item Некоторые tone-mapping кривые \textbf{уже включают} в себя гамма-коррекцию!
\end{itemize}
\end{frame}

\begin{frame}[fragile]
\frametitle{Коррекция гаммы: ссылки}
\begin{itemize}
\item \href{https://en.wikipedia.org/wiki/Gamma_correction}{\texttt{en.wikipedia.org/wiki/Gamma\_correction}}
\item \href{http://blog.johnnovak.net/2016/09/21/what-every-coder-should-know-about-gamma}{\texttt{What every coder should know about gamma}}
\item \href{http://filmicworlds.com/blog/linear-space-lighting-i-e-gamma}{\texttt{Linear-space lighting (i.e. gamma)}}
\item \href{https://learnopengl.com/Advanced-Lighting/Gamma-Correction}{\texttt{Learnopengl.com tutorial}}
\end{itemize}
\end{frame}

\begin{frame}[fragile]
\frametitle{sRGB}
\begin{itemize}
\item По тем же причинам, по которым существует нелинейная гамма у мониторов (отдать больше бит под тёмные цвета), хочется хранить текстуры не в сыром формате, а в гамма-преобразованном: значение пикселя \begin{math}V\end{math} зависит от желаемой интенсивности \begin{math}I\end{math} как \begin{math}V \sim I^\frac{1}{\gamma}\end{math}
\pause
\item Такой формат хранения цвета называется sRGB
\begin{itemize}
\item N.B. Точная формула преобразования sRGB отличается от \begin{math}I^\frac{1}{2.2}\end{math}, но очень близка к ней
\end{itemize}
\pause
\item Почти все изображения хранятся именно в таком формате!
\pause
\item Проверить формат можно любым редактором изображений или программой \mintinline{cpp}|identify| пакета \mintinline{cpp}|imagemagick|
\end{itemize}
\end{frame}

\begin{frame}[fragile]
\frametitle{sRGB}
\begin{itemize}
\item При чтении из sRGB-текстуры нужно возводить прочитанное значение в степень 2.2
\pause
\item При записи в sRGB-текстуру нужно возводить записываемое значение в степень 1/2.2
\pause
\item Если нам нужно просто вывести картинку на экран, можно игнорировать оба эффекта -- они гасят друг друга
\pause
\item В WebGPU есть встроенная поддержка sRGB
\end{itemize}
\end{frame}

\begin{frame}[fragile]
\frametitle{sRGB в WebGPU}
\begin{itemize}
\item У некоторых форматов текстур есть sRGB-вариант: например, \mintinline{rust}|Rgba8Unorm| \(\rightarrow\) \mintinline{rust}|Rgba8UnormSrgb|
\pause
\item Нельзя создать текстуру с sRGB-форматом, но можно создать texture view с соответствующим sRGB-форматом (тогда этот формат нужно также указать в \mintinline{rust}|view_formats| при создании текстуры)
\pause
\item Чтение из такого texture view автоматически делает преобразование из sRGB в \textit{линейный} цвет (т.е. в физическую интенсивность света)
\pause
\item Рендеринг в color attachment с таким texture view автоматически делает преобразование \textit{линейного} цвета в sRGB
\pause
\item В конструкторе \mintinline{rust}|wgpu_app| в практиках есть параметр \mintinline{rust}|srgb| -- он включает sRGB-формат для нашего экрана
\end{itemize}
\end{frame}

\begin{frame}[fragile]
\frametitle{sRGB mipmaps}
\begin{itemize}
\item Генерировать mipmaps для sRGB-текстур нужно с учётом sRGB-преобразования
\pause
\item Проще всего воспользоваться sRGB-форматами texture view, чтобы преобразование было автоматическим
\pause
\item Тогда можно переиспользовать тот же шейдер, что и для не-sRGB mipmaps, но придётся создать отдельный render pipeline (он зависит от формата color attachment)
\end{itemize}
\end{frame}

\begin{frame}[fragile]
\frametitle{sRGB-correct mipmap}
\slideimage{srgb-mipmap.png}
\end{frame}

\begin{frame}[fragile]
\frametitle{sRGB: ссылки}
\begin{itemize}
\item \href{https://en.wikipedia.org/wiki/SRGB}{\texttt{en.wikipedia.org/wiki/SRGB}}
\item \href{https://stackoverflow.com/a/12894053/2315602}{\texttt{Stackoverflow post on sRGB}}
\item \href{https://www.khronos.org/opengl/wiki/Framebuffer#Colorspace}{\texttt{sRGB framebuffers}}
\end{itemize}
\end{frame}

\begin{frame}[fragile]
\frametitle{Color banding}
\slideimage{banding.jpg}
\end{frame}

\begin{frame}[fragile]
\frametitle{Color banding}
\begin{itemize}
\item Артефакт, когда чётко видна граница между областями, заполненными одним цветом
\pause
\item Обычно не сильно заметен, особенно после наложения текстур и пост-обработки
\pause
\item Хорошо заметен на монотонных поверхностях (напр. небо)
\pause
\item Сильнее заметен с тёмными цветами (из-за особенностей человеческого зрения)
\pause
\item Может проявляться из-за ограниченности вычислений (напр. маленькое число сэмплов при SSAO)
\pause
\item Может усиливаться при неаккуратной работе с форматами хранения и нелинейными преобразованиями
\end{itemize}
\end{frame}

\begin{frame}[fragile]
\frametitle{Color banding: пример возникновения}
\begin{itemize}
\item Рисуем сцену, записываем линейную интенсивность в текстуру в формате \mintinline{cpp}|Rgba8Unorm|
\pause
\item Выводим сцену на экран с гамма-коррекцией: значения из текстуры возводятся в степень 1/2.2
\pause
\item Значение текстуры 0 переходит в значение 0 на экране
\pause
\item Значение текстуры 1 переходит в значение \begin{math}255 \cdot \left(\frac{1}{255}\right)^\frac{1}{2.2} = 21\end{math} на экране
\pause
\item Колоссальная потеря точности: значения от 1 до 20 не используются!
\pause
\item Для избавления можно увеличить битность промежуточной текстуры до \mintinline{cpp}|Rgba16Unorm|
\end{itemize}
\end{frame}

\begin{frame}[fragile]
\frametitle{Color banding}
\begin{itemize}
\item Banding, связанный с потерей точности, обычно решают увеличением точности (8-bit \begin{math}\longrightarrow\end{math} 16-bit)
\pause
\item Banding, связанный ограниченными вычислениями, обычно решают рандомизацией, разменивая banding на шум
\pause
\item Banding, связанный с ограниченной точностью самого экрана, обычно решают с помощью \textit{дизеринга}
\end{itemize}
\end{frame}

\begin{frame}[fragile]
\frametitle{Dithering}
\slideimage{dither.png}
\end{frame}

\begin{frame}[fragile]
\frametitle{Dithering}
\begin{itemize}
\item Способ рисования полутонов за счёт варьирования количества пикселей двух цветов
\pause
\item Работает только для больших (по размеру) градиентов
\pause
\item Работает при любой точности формата хранения, визуально увеличивая её
\pause
\item Может использоваться для рисования полутонов между областями, окрашенными в соседние представимые цвета
\pause
\item Может использоваться как художественный приём (имитация пиксель-арт графики, имитация очень старых игр)
\end{itemize}
\end{frame}

\begin{frame}[fragile]
\frametitle{Dithering}
\slideimage{dither2.png}
\end{frame}

\begin{frame}<1>[fragile,label=dither]
\frametitle{Dithering: реализация}
\begin{itemize}
\item Допустим, мы рисуем в однобитное (чёрно-белое) изображение, и хотим добиться серого цвета с помощью дизеринга
\pause
\item Нам нужен способ определить, какой из пикселей закрасить чёрным, а какой -- белым
\pause
\item При одноразовой обработке изображения можно оптимизировать паттерн дизеринга под конкретное изображение
\pause
\item В real-time графике это обычно делается с помощью т.н. dither mask / dither matrix
\end{itemize}
\end{frame}

\begin{frame}[fragile]
\frametitle{Dithering}
\slideimage{dither3.png}
\end{frame}

\againframe<1->{dither}

\begin{frame}[fragile]
\frametitle{Dithering: реализация}
\begin{itemize}
\item Dither mask -- одноканальная текстура, хранящая для каждого пикселя пороговое значение, начиная с которого его надо рисовать белым
\pause
\item Например, для серого цвета с яркостью 0.25 мы нарисуем белыми те пиксели, для которых значение в dither map \begin{math}\leq 0.25\end{math}
\pause
\item Например, для серого цвета с яркостью 0.5 мы нарисуем белыми те пиксели, для которых значение в dither map \begin{math}\leq 0.5\end{math}
\pause
\item Например, для серого цвета с яркостью 0.75 мы нарисуем белыми те пиксели, для которых значение в dither map \begin{math}\leq 0.75\end{math}
\end{itemize}
\end{frame}

\begin{frame}[fragile]
\frametitle{Dither mask}
\slideimage{dither-mask.png}
\end{frame}

\begin{frame}[fragile]
\frametitle{Dithering: реализация}
\begin{itemize}
\item Для дизеринга при рисовании в 8-битный буфер:
\begin{itemize}
\item Хотим записать значение \begin{math}v \in [0, 1]\end{math}
\item Вычисляем \begin{math}v' = 255 \cdot v\end{math}
\item Вычисляем \begin{math}v_0 = floor(v')\end{math}
\item Вычисляем \begin{math}v_1 = v_0 + 1\end{math}
\item Вычисляем \begin{math}dv = v' - v0\end{math}
\item Если значение в dither mask меньше \begin{math}dv\end{math}, рисуем значение \begin{math}v_1 / 255\end{math}, иначе \begin{math}v_0 / 255\end{math}
\end{itemize}
\pause
\item Дизеринг нужно делать для каждого канала (RGB) по отдельности, \textbf{после} применения tone mapping и гамма-коррекции
\pause
\item Обычно dither mask имеет маленький размер (напр. 16x16) и циклически (\mintinline{cpp}|addressMode = Repeat|) повторяется на весь экран; её пиксели пробегают все возможные значения от 0 до 255
\end{itemize}
\end{frame}

\begin{frame}[fragile]
\frametitle{Dithering: реализация}
\usemintedstyle{lightbulb}
\begin{tcblisting}{listing engine=minted, minted language=wgsl, minted options={fontsize=\scriptsize}, listing only, colback=codebg, colframe=codebg, rounded corners}
var color = ...;

let c = color * 255.0;
let c0 = floor(c);
let c1 = c0 + vec3f(1.0);
let dc = c - c0;

let threshold = texture(dither_mask, pixel_coord).r;

color = select(c0, c1, dc > threshold) / 255.0;
\end{tcblisting}
\end{frame}

\begin{frame}[fragile]
\frametitle{Генерация dither mask}
\begin{itemize}
\item Bayer (ordered) dithering -- очень быстро генерируется (явная формула), могут быть заметны артефакты из-за регулярности паттерна
\pause
\item Blue noise dithering -- использует синий шум (напр. алгоритм void-and-cluster) для генерации маски, лучше распределение пикселей, нет артефактов из-за регулярности
\pause
\item Готовые текстуры с blue noise dither mask разных размеров можно легко найти онлайн
\end{itemize}
\end{frame}

\begin{frame}[fragile]
\frametitle{Сравнение ordered dither и blue noise dither}
\slideimage{dither-mask-compare.png}
\end{frame}

\begin{frame}[fragile]
\frametitle{Screen-door transparency}
\begin{itemize}
\item Дизеринг часто применяют как дешёвый аналог прозрачности
\pause
\item Если значение в альфа-канале больше значения из dither mask, то рисуем пиксель, иначе -- не рисуем (\mintinline{wgsl}|discard|)
\pause
\item Не требует сортировки объектов по расстоянию, работает с deferred shading'ом
\pause
\item Выглядит заметно хуже, часто применяется для временно прозрачных объектов (появляющихся/исчезающих), далёких объектов (деревья, кусты), или улучшается пост-обработкой кадра
\end{itemize}
\end{frame}

\begin{frame}[fragile]
\frametitle{Screen-door transparency}
\slideimage{dither-transparency.png}
\end{frame}

\begin{frame}[fragile]
\frametitle{Return of the Obra Dinn (2018)}
\slideimage{obra-dinn.png}
\end{frame}

\begin{frame}[fragile]
\frametitle{Dithering: ссылки}
\begin{itemize}
\item \href{https://en.wikipedia.org/wiki/Dither}{\nolinkurl{en.wikipedia.org/wiki/Dither}}
\item \href{https://surma.dev/things/ditherpunk}{\nolinkurl{surma.dev/things/ditherpunk}}
\item \href{http://alex-charlton.com/posts/Dithering_on_the_GPU}{\nolinkurl{alex-charlton.com/posts/Dithering_on_the_GPU}}
\item \href{https://forums.tigsource.com/index.php?topic=40832.msg1363742}{\texttt{Dithering in Return of the Obra Dinn}}
\item \href{https://www.shadertoy.com/view/Dscczs}{\texttt{Shadertoy с примером дизеринга}}
\item \href{https://momentsingraphics.de/BlueNoise.html}{\texttt{Free blue noise textures}}
\end{itemize}
\end{frame}

\end{document}